\documentclass[11pt,a4paper]{article}

% ---------------------------------------------------------
% PACKAGES
% ---------------------------------------------------------
\usepackage[utf8]{inputenc}
\usepackage[T1]{fontenc}
\usepackage{lmodern}
\usepackage{amsmath, amssymb, amsfonts, bm}
\usepackage{geometry}
\usepackage{hyperref}
\usepackage{physics}
\usepackage{graphicx}
\usepackage{cite}
\usepackage{tensor}
\usepackage{mathtools}
\usepackage{bbm}

\geometry{margin=2.5cm}

% ---------------------------------------------------------
% TITLE
% ---------------------------------------------------------
\title{\textbf{Companion I (v49.1)}\\[0.2cm]
\textbf{Global CPT Symmetry and Structural Gravity}\\[0.3cm]
\large Foundations of Relaxation-Driven Cyclic Cosmology (RDCC)}

\author{Michael Lehmann\\
Independent Researcher\\
\texttt{mi.lehmann@gmx.de}}

\date{September 2026}

\begin{document}
\maketitle

% ---------------------------------------------------------
% ABSTRACT
% ---------------------------------------------------------
\begin{abstract}
Relaxation-Driven Cyclic Cosmology (RDCC) describes the Universe as a globally
CPT-symmetric quantum state defined on a bipartite Hilbert space
$\mathcal{H} = \mathcal{H}_{(+)} \otimes \mathcal{H}_{(-)}$. The visible and
mirror sectors arise as reduced density matrices of a single global state and
possess thermodynamic arrows of time that point in opposite directions. Gravity
is the only interaction that couples to the \emph{total} stress-energy tensor
and therefore probes the full global state. In v49.1, the foundational
structure of RDCC is fully Phase-1-integrated: the inter-sector operator
$T^{(+)}_{\mu\nu}T^{(-)\mu\nu}$ is consistently treated as a dimension-8
operator, the relaxation parameter $\alpha(t)$ is globally defined and
RG-stabilized, and the mirror sector forms diffuse, non-compact halos. This
Companion establishes the conceptual and mathematical backbone for all
subsequent RDCC components.
\end{abstract}

\tableofcontents
\newpage
% ============================================================
\section{Introduction}
% ============================================================

Standard cosmology assumes a single thermodynamic sector evolving forward in
time. RDCC extends this picture by proposing that the Universe is fundamentally
a \emph{global} CPT-symmetric quantum state defined on a bipartite Hilbert
space,
\begin{equation}
\mathcal{H} = \mathcal{H}_{(+)} \otimes \mathcal{H}_{(-)}.
\end{equation}

Sectoral observers access only one factor of this Hilbert space and therefore
perceive an effectively time-oriented, thermodynamic Universe. The mirror
sector is the CPT-conjugate projection of the visible sector and possesses a
thermodynamic arrow of time that points in the opposite direction.

Gauge interactions act strictly within one sector. Gravity is distinguished by
its coupling to the \emph{total} stress-energy tensor,
\begin{equation}
T^{\mathrm{tot}}_{\mu\nu}
=
T^{(+)}_{\mu\nu}
+
T^{(-)}_{\mu\nu},
\end{equation}
and therefore probes the full global state. This structural asymmetry is the
origin of the gravitational contribution of the mirror sector to the visible
sector.

In v49.1, the foundational structure of RDCC is refined and fully Phase-1-
integrated: the inter-sector operator is treated as a dimension-8 operator, the
relaxation parameter $\alpha(t)$ is globally defined and RG-stabilized, and the
mirror sector forms diffuse, non-compact halos.
% ============================================================
\section[Global CPT Structure]{Global CPT Structure and Bipartite Hilbert Space}
% ============================================================

The global state may be written in Schmidt form,
\begin{equation}
\ket{\Psi}
=
\sum_{i} \lambda_{i}\,
\ket{i}_{(+)} \otimes \ket{i}_{(-)},
\end{equation}
with coefficients $\lambda_{i}$ encoding inter-sector correlations.

A visible-sector observer perceives the reduced density matrix,
\begin{equation}
\rho_{(+)}
=
\mathrm{Tr}_{(-)}\left(\ket{\Psi}\bra{\Psi}\right),
\end{equation}
and all local observables are defined on $\mathcal{H}_{(+)}$ alone.

The mirror sector arises analogously,
\begin{equation}
\rho_{(-)}
=
\mathrm{Tr}_{(+)}\left(\ket{\Psi}\bra{\Psi}\right).
\end{equation}

Both reduced states are mixed even when the global state is pure. This
decoherence is structural: it arises from tracing out the degrees of freedom of
the opposite sector.

Global CPT symmetry requires
\begin{equation}
\mathrm{CPT}\ket{\Psi} = \ket{\Psi},
\end{equation}
and enforces opposite thermodynamic arrows of time,
\begin{equation}
t_{(-)} = - t_{(+)}.
\end{equation}
% ============================================================
\section{Structural Gravity and Sectoral Projection}
% ============================================================

Gravity responds to the full global state. The global Einstein equation is
\begin{equation}
G_{\mu\nu}[g]
=
8\pi G_{N}\left(
T^{(+)}_{\mu\nu}
+
T^{(-)}_{\mu\nu}
\right).
\label{eq:globalEinstein}
\end{equation}

A sectoral observer infers the projected Einstein equation,
\begin{equation}
G_{\mu\nu}[g]
=
8\pi G_{\mathrm{eff}}\,T^{(+)}_{\mu\nu},
\label{eq:sectoralEinstein}
\end{equation}
with an effective gravitational coupling $G_{\mathrm{eff}}$.

Linearizing the global response yields
\begin{equation}
\delta G_{\mu\nu}
=
8\pi G_{N}\left(
\delta T^{(+)}_{\mu\nu}
+
\delta T^{(-)}_{\mu\nu}
\right).
\end{equation}

Equating global and sectoral responses gives
\begin{equation}
G_{\mathrm{eff}}
=
G_{N}\left(
1
+
\frac{\delta T^{(-)}_{\mu\nu}}{\delta T^{(+)}_{\mu\nu}}
\right).
\end{equation}

For homogeneous FRW backgrounds,
\begin{equation}
\frac{\delta T^{(-)}_{\mu\nu}}{\delta T^{(+)}_{\mu\nu}}
=
\frac{\rho_{(-)}}{\rho_{(+)}},
\end{equation}
leading to
\begin{equation}
G_{\mathrm{eff}}
=
G_{N}\left(
1
+
\frac{\rho_{(-)}}{\rho_{(+)}}
\right).
\label{eq:GeffFinal}
\end{equation}
% ============================================================
\section[Dimension-8 Operator and Global alpha(t)]{Dimension-8 Operator and Global $\alpha(t)$}
% ============================================================

The inter-sector coupling is encoded in the dimension-8 operator
\begin{equation}
\mathcal{O}_{\mathrm{int}}
=
T^{(+)}_{\mu\nu}T^{(-)\mu\nu},
\end{equation}
with effective action
\begin{equation}
S_{\mathrm{int}}
=
\int d^4x \sqrt{-g}\,
\frac{1}{M^{4}}\,
T^{(+)}_{\mu\nu}T^{(-)\mu\nu}.
\end{equation}

The relaxation parameter $\alpha(t)$ is defined globally as the infrared
response of the metric to this operator,
\begin{equation}
\alpha(t) \propto \frac{1}{M^{4}}.
\end{equation}

RG flow stabilizes the coupling in the infrared,
\begin{equation}
\dot{\alpha}(t) = -\Gamma_{\mathrm{IR}}\,\alpha(t),
\end{equation}
with solution
\begin{equation}
\alpha(t)
=
\alpha_{\mathrm{IR}}
+
\left(\alpha_{0}-\alpha_{\mathrm{IR}}\right)
e^{-\Gamma_{\mathrm{IR}} t}.
\end{equation}
% ============================================================
\section{Bounce Structure}
% ============================================================

At the bounce,
\begin{equation}
H(t_{b}) = 0,
\qquad
\alpha(t_{b}) = 0,
\end{equation}
and the two sectors become maximally correlated. The global state is perfectly
CPT-symmetric, and the metric responds to the total stress-energy tensor.

The visible sector perceives a nonsingular bounce because the mirror-sector
energy density appears as a purely gravitational contribution and provides the
structural origin of dark matter.

% ============================================================
\section{Conclusion}
% ============================================================

This Companion develops the refined foundational structure of RDCC in its
v49.1 form. Gravity acts as the global bridge between the CPT-conjugate
sectors, responding to the sum of their stress-energy tensors. The inter-sector
operator is consistently treated as a dimension-8 operator, restoring EFT
coherence and providing the microscopic origin of the global relaxation
parameter $\alpha(t)$.

These foundations support the relaxation sector (Companion III), the tensor
sector (Companion II and XIV), the thermodynamic sector (Companion IV), the
ultraviolet structure (Companion X), and the survey interface (Companion XII).

\end{document}
